\section{Introduction}

The Dirichlet beta function is
\[
 \beta(s)=\sum_{m=0}^{\infty}\frac{(-1)^m}{(2m+1)^s}.
\]
Its values at positive odd integers are nonzero rational multiples of
powers of $\pi$. The arithmetic of its even values is much less understood.
One way to measure progress without settling an individual value is to
bound the dimension
\begin{equation}\label{eq:dimension}
 D(s)=\dim_{\Q}\Span_{\Q}\{\beta(2),\beta(4),\ldots,\beta(s)\},
 \qquad s\geq2\ \text{even}.
\end{equation}
A lower bound $D(s)\geq d$ means that at least $d$ values in this finite
list are linearly independent over the rationals. It implies that at
least $d-1$ of them are irrational, but does not identify those values.

Rivoal and Zudilin \cite[Theorem~1]{RZ03} proved a logarithmic lower bound
with leading constant $1/(2+\log2)$ for the span after adjoining $1$.
Fischler \cite[Theorem~1]{F20} subsequently obtained
\begin{equation}\label{eq:oldbound}
 D(s)\geq\frac{1+o(1)}{1+\log2}\log s
 \qquad(s\longrightarrow\infty,\ s\text{ even}),
\end{equation}
as a special case of his result for periodic Dirichlet series.
His later use of Siegel's lemma to construct rational functions with
small coefficients gives a bound of order $\sqrt{s/\log s}$ for odd
zeta values \cite{F26}.

Mistiaen \cite[Theorem~1.0.5 and Remark~3]{M25} extends that method to
Dirichlet characters whose conductors are not divisible by $4$. The
restriction excludes the character of conductor $4$ associated with
$\beta$, so that theorem does not cover the even beta values.
We obtain the following bound for the beta function.

\begin{theorem}\label{thm:main}
For every sufficiently large even integer $s$,
\[
 D(s)\geq0.21\sqrt{\frac{s}{\log s}}.
\]
\end{theorem}

The theorem concerns linear independence in a growing family. It does
not, in particular, establish the irrationality of $\beta(4)$.
A recent preprint of Kumar and Tosi \cite{KT26} claims a bound of
square-root order for the \emph{number of irrational values} of real
Dirichlet $L$-functions. The proof contains unresolved arithmetic
errors,\footnote{For example, in the notation of \cite[Section~3.3]{KT26},
take $a=b=1$ and $h=y=\ell=0$. For sufficiently large admissible $n$,
$K_0(0)=\prod_{j=0}^{(2r+T)n}(1-rn+j)=0$, because the factor
$j=rn-1$ vanishes, whereas $g_{0,1,1}(0)=1$ by the empty-product
convention. This contradicts the finite $p_n$-adic valuation asserted
in Lemma~18. This objection concerns the proof as written and does not
disprove the main theorem.}
so we do not regard that version as establishing the claimed bound.
Our proof is independent of theirs.
Theorem~\ref{thm:main} concerns the stronger
quantity \eqref{eq:dimension}: irrational numbers can still satisfy many
rational linear relations.

Specializing general results on polylogarithms with periodic coefficients
\cite{M26} gives bounds involving both even and odd beta values.
However, the odd values $\beta(1),\beta(3),\ldots$ are already linearly
independent over the rationals: they are nonzero rational multiples of
distinct powers of $\pi$, and $\pi$ is transcendental.
The lower bound supplied by those results can therefore be met by the odd
values alone, without establishing the desired independence among the even
values. Our construction avoids this issue by producing linear forms
involving only $1$ and the even beta values.

\subsection{How the proof works}

For a fixed even cutoff $s=2v$, the proof constructs integer linear forms
in $1$ and the even beta values:
\[
 \lambda_0+\lambda_1\beta(2)+\cdots+\lambda_v\beta(2v),
 \qquad \lambda_0,\ldots,\lambda_v\in\Z.
\]
If all these beta values were rational, any nonzero such form would have
absolute value bounded below by the reciprocal of a common denominator.
Linear-independence criteria refine this observation to give a lower bound
on the dimension of their span.

With this set of beta values fixed, we use a criterion of Fischler, stated in
Proposition~\ref{prop:criterion}. It requires a family of forms at each
large index $n$, with controlled growth of their integer coefficients and
exponential decay of their values. It also requires a rank condition on
the coefficient matrix, which prevents the whole family from vanishing.
Writing these bounds as $B^{n+o(n)}$ for the coefficients and
$\alpha^{n+o(n)}$ for the values, with $B>1$ and $0<\alpha<1$, the
criterion gives $D(s)\geq-\log\alpha/\log B$ after removing the
adjoined number $1$. Thus constructing a small expression is only the first step:
we must control its size after clearing denominators and establish the
required rank condition.

The starting point is Fischler's rational function
\[
 F_n(t)=\sum_{i=1}^{a}\sum_{j=0}^{n}\frac{c_{i,j}}{(t+j)^i},
 \qquad c_{i,j}\in\Z,
\]
whose coefficients have controlled size and whose many initial zero
Laurent coefficients make it decay rapidly far from its poles
\cite[Theorem~4]{F26}. The coefficients are supplied by an existence
theorem; no explicit formula for them is needed.
We sum this rational function and its first $h$ derivatives over
half-integers. These give $h+1$ families of forms. The additional
families provide the independent functions needed in the rank argument;
their cost is to raise the largest possible beta argument from $a$ to
$a+h$.

Half-integer arguments already occur in classical approximations to beta
values, including the hypergeometric construction of
Rivoal and Zudilin \cite[Section~2]{RZ03}.
They are natural here because powers of $t+\tfrac12$ are proportional to
powers of the odd integers $2t+1$ in the beta series. Alternating sums
therefore produce beta values. Combining the positive and negative
half-integer sums removes the odd beta values.
Section~\ref{sec:forms} carries this out using generating functions
evaluated at $z=-1$.

These generating functions satisfy a rational differential system whose
coefficients have no pole at $-1$: it is an \emph{ordinary point} of the
system. This avoids the singular evaluation point identified in
\cite[Section~2]{M25}. At an ordinary point, the initial values determine
an analytic solution, and its Taylor coefficients express derivatives
in terms of those initial values. This is useful for the arithmetic of
our linear forms. The sums first produce finite correction terms with
rational coefficients. We remove those terms by differentiation
\emph{before} clearing denominators. Section~\ref{sec:arithmetic} then
bounds the denominators through the local differential system, and
Section~\ref{sec:estimates} compares the resulting coefficient growth
and decay.

The remaining task is the rank condition. Section~\ref{sec:rank} follows
Fischler's multiplicity method \cite[Section~4.5]{F26}: failure of the
rank condition would force certain polynomial combinations of solutions
to have more vanishing than their degrees permit. Each of the $h+1$
families contributes one such combination, and the final count requires
$h$ to grow proportionally to $a$. Our system has an extra diagonal
term, so we first make the change $z=w^2$ and rescale
some components of the solution. This removes that term and allows the
independence argument to work. The substitution doubles the relevant
polynomial degrees, but the evaluation point now has two preimages,
$\ii$ and $-\ii$. Counting the vanishing at both recovers the required
inequality.

The size of the final bound comes from balancing two costs. The
rational-function construction has a parameter $\Omega$ that improves
the decay exponent by a term proportional to $\Omega$, but increases
the coefficient-growth exponent by a term of order $\Omega^2/a$.
With the other auxiliary parameters fixed and $h$ proportional to $a$,
the denominator estimate contributes only a term of order $\log a$.
Choosing $\Omega$ of order
$\sqrt{a\log a}$ therefore keeps $\log B$ of order $\log a$ and
makes $-\log\alpha$ of order $\sqrt{a\log a}$.
Their quotient has order $\sqrt{a/\log a}$; the largest beta argument
$a+h$ is proportional to $a$. Section~\ref{sec:conclusion} makes this
balance precise and uses Fischler's numerical parameter choice to
obtain the constant $0.21$.

Section~\ref{sec:inputs} states the criterion and the rational-function
input precisely. The multiplicity estimate is stated in
Section~\ref{subsec:multiplicity}, alongside the rank argument that uses it.

\subsection{Conventions}

All logarithms are natural. We write
$(i)_p=i(i+1)\cdots(i+p-1)$, with $(i)_0=1$, and
$\Ht_d=\sum_{j=1}^d1/j$ for the $d$th harmonic number.
The symbol $\ii$ denotes the imaginary unit.
The auxiliary parameters of the construction are fixed while $n$ tends
to infinity through an indicated infinite set of integers.
Every $o(n)$ in that stage of the proof is uniform in the finite indices
and in the specified derivative range. Only in the last section do we
let the parameter controlling the largest beta argument tend to infinity.

